% =====================================================================
% Paper 38 (standalone): SU(2)/S^3 Latremoliere propinquity convergence
% on Camporesi-Higuchi truncated metric spectral triples
%
% J. Geom. Phys. style; companion to Leimbach--van Suijlekom 2024
% (Adv. Math. 439, 109496) for flat tori.
%
% Status: first draft, May 2026, post-WH1-PROVEN closure (2026-05-06).
% Built from internal proof memos: r25_l2_proof_memo.md,
% r25_l2_quantitative_rate_memo.md, r25_l3_proof_memo.md,
% r25_l4_proof_memo.md, r25_l5_proof_memo.md, wh1_r35_full_dirac_memo.md.
% =====================================================================
\documentclass[11pt,a4paper]{article}

\usepackage[T1]{fontenc}
\usepackage[utf8]{inputenc}
\usepackage{amsmath, amssymb, amsthm, mathrsfs, mathtools}
\usepackage{geometry}
\geometry{margin=1in}
\usepackage{hyperref}
\hypersetup{colorlinks=true, linkcolor=blue!50!black, citecolor=blue!50!black, urlcolor=blue!50!black}
\usepackage{microtype}
\usepackage{graphicx}
\usepackage{booktabs}
\usepackage[numbers,sort&compress]{natbib}
\usepackage{xcolor}

\theoremstyle{plain}
\newtheorem{theorem}{Theorem}[section]
\newtheorem{lemma}[theorem]{Lemma}
\newtheorem{proposition}[theorem]{Proposition}
\newtheorem{corollary}[theorem]{Corollary}
\theoremstyle{definition}
\newtheorem{definition}[theorem]{Definition}
\newtheorem{remark}[theorem]{Remark}
\newtheorem{example}[theorem]{Example}

\newcommand{\nmax}{n_{\max}}
\newcommand{\sthree}{S^{3}}
\newcommand{\Hilb}{\mathcal{H}}
\newcommand{\Op}{\mathcal{O}}
\newcommand{\Tcal}{\mathcal{T}}
\newcommand{\Acal}{\mathcal{A}}
\newcommand{\opnorm}[1]{\lVert #1 \rVert_{\mathrm{op}}}
\newcommand{\Lipnorm}[1]{\lVert #1 \rVert_{\mathrm{Lip}}}
\newcommand{\norm}[1]{\lVert #1 \rVert}
\newcommand{\R}{\mathbb{R}}
\newcommand{\C}{\mathbb{C}}
\newcommand{\T}{\mathbb{T}}
\newcommand{\N}{\mathbb{N}}
\newcommand{\Z}{\mathbb{Z}}
\newcommand{\Q}{\mathbb{Q}}
\newcommand{\SU}{\mathrm{SU}}
\newcommand{\SO}{\mathrm{SO}}
\newcommand{\Vol}{\mathrm{Vol}}
\newcommand{\Tr}{\mathrm{Tr}}
\newcommand{\Mat}{M}
\newcommand{\DCH}{D_{\mathrm{CH}}}
\newcommand{\KS}{K}
\newcommand{\Khat}{\widehat{K}}

\title{Latr\'emoli\`ere propinquity convergence of \\
truncated Camporesi--Higuchi spectral triples on $\sthree$}
\author{J.~Loutey\thanks{Independent researcher. \texttt{jloutey@gmail.com}.
This work was developed in an AI-augmented agentic workflow with
Anthropic's Claude. Implementation, internal memos, and bibliographic
collation were performed collaboratively; all theorems, design choices,
and editorial decisions are the author's responsibility.}}
\date{May 2026 (preprint)}

\begin{document}
\maketitle

\begin{abstract}
We prove that the Connes--van~Suijlekom truncations of the round
$\sthree = \SU(2)$ Camporesi--Higuchi spectral triple converge to the
continuum spectral triple in the Latr\'emoli\`ere quantum
Gromov--Hausdorff propinquity, with explicit asymptotic rate
\(
\Lambda(\Tcal_{\nmax}, \Tcal_{\sthree}) = O\!\bigl(\log \nmax / \nmax\bigr)
\)
and asymptotic constant $4/\pi$. The proof proceeds via five lemmas:
(L1$'$) construction of a truncated operator system on the full
chirality-doubled spinor bundle with finite Connes distance on every
cross-shell pair; (L2) construction of a positive central spectral
Fej\'er kernel on $\SU(2)$ with closed-form mass-concentration moments
and Stein--Weiss-type quantitative rate $\gamma_n \sim (4/\pi)\log n / n$;
(L3) a Lipschitz comparison bound
$\opnorm{[\DCH, \Mat_f]} \leq C_3\,\norm{\nabla f}_{L^\infty}$ with
$C_3 = 1$; (L4) a Berezin reconstruction map that is positive,
contractive, an approximate identity, and L3-compatible; (L5)
mechanical assembly of the propinquity bound via a tunneling pair.
The asymptotic constant $4/\pi = \Vol(S^2)/\pi^2$ identifies as the
Hopf-base measure factor in the standard $\SU(2)/U(1)$ Haar
normalisation, exhibiting the expected one-log slowdown of the
non-abelian SU(2) rate relative to the
Leimbach--van~Suijlekom~\cite{leimbach_vs2024} torus rate $O(1/\nmax)$.
The result extends Connes--van~Suijlekom~\cite{connes_vs2021} and the
subsequent flat-structure literature
(\cite{hekkelman2022, hekkelman_mcdonald2024, leimbach_vs2024}) to the
first physically natural compact non-abelian case, and supplies the
metric foundation for almost-commutative spectral-triple constructions
hosted on $\sthree$~\cite{marcolli_vs2014}.
\end{abstract}

\bigskip\noindent
\textbf{MSC2020 classifications:}\ 58B34 (noncommutative geometry),
46L87 (noncommutative Riemannian and metric geometry),
46L52 (noncommutative function spaces),
22E45 (representations of compact Lie groups).

\medskip\noindent
\textbf{Keywords:}\ spectral triple, quantum Gromov--Hausdorff
propinquity, Camporesi--Higuchi Dirac, operator system, Berezin
reconstruction, Peter--Weyl, $\SU(2)$.

% =====================================================================
\section{Introduction}\label{sec:intro}
% =====================================================================

A central problem in noncommutative metric geometry is whether
finite-dimensional truncations of a Riemannian spectral triple
converge, in a suitable quantum Gromov--Hausdorff sense, to the
continuum geometry. Connes and van~Suijlekom~\cite{connes_vs2021}
introduced the framework of \emph{truncated operator systems}
$\Op_\Lambda = P_\Lambda \mathcal{A} P_\Lambda$ obtained by spectral
projection onto the modes of $D$ with $|D| \le \Lambda$, and showed
that the natural propagation invariant $\mathrm{prop}(\Op_\Lambda) = 2$
is realised on the Toeplitz operator system on $S^1$. They left
explicit convergence to the continuum, in the Latr\'emoli\`ere
propinquity~\cite{latremoliere2018, latremoliere2016} or any
quantitative quantum Gromov--Hausdorff distance, as an open question,
deferring it to ``elsewhere'' on three separate occasions in the
text\footnote{See \cite[lines~220--221, 1101, 2616--2618 of
\texttt{arXiv:2004.14115v2}]{connes_vs2021}.}.

The flat-structure case has been worked out in several subsequent
papers. Hekkelman~\cite{hekkelman2022} treated the circle $S^1$
with the Toeplitz/Fej\'er kernel reconstruction, and
Hekkelman--McDonald~\cite{hekkelman_mcdonald2024} extended the
analysis to $T^d$. Leimbach and van~Suijlekom~\cite{leimbach_vs2024}
proved propinquity convergence of truncated tori with explicit rate
$O(1/\Lambda)$, using the standard
Schur-multiplier characterisation of the abelian-compact-group Fej\'er
mechanism. A complementary line of work due to
Hekkelman--McDonald--van~Suijlekom~\cite{ucp_maps_2024} treats
\emph{Berezin--Toeplitz} reconstruction on $S^2$, using the K\"ahler
structure of the coadjoint orbit. All of these results address flat
abelian or K\"ahler structures.

Adjacent activity in the broader propinquity program addresses
orthogonal corners of the same general framework.
Toyota~\cite{toyota2023} establishes quantum-Gromov--Hausdorff
convergence of spectral truncations for discrete groups of
polynomial growth, a scope which excludes compact Lie groups by
definition. Farsi--Latr\'emoli\`ere~\cite{farsi_latremoliere2024}
prove spectral continuity for collapse of products of spectral
triples containing at least one \emph{abelian} factor, an
abelian-factor restriction which the present paper lifts.
Farsi--Latr\'emoli\`ere~\cite{farsi_latremoliere2025} establish
continuity of the spectral propinquity along analytic paths of
Riemannian metrics on a fixed closed spin manifold, a
continuous-family setup distinct from the discrete-truncation
sequence considered here. Hekkelman--McDonald~\cite{hekkelman_mcdonald2024b}
develop a noncommutative integral and a Szeg\"o limit formula for
spectrally truncated triples in the Connes--van~Suijlekom paradigm,
a companion direction complementary to the metric-side convergence
proved here.

The most natural physically motivated non-K\"ahler compact case ---
the round $\sthree = \SU(2)$ with a Dirac spectral triple --- has
remained open.

The present paper closes that gap. Let $\Tcal_{\sthree} =
(C^\infty(\sthree), L^2(\sthree, \Sigma), \DCH)$ denote the
\emph{Camporesi--Higuchi spectral triple}, with $\DCH$ the Dirac
operator on the unit round $\sthree$ in the spinor harmonic basis
of Camporesi--Higuchi~\cite{camporesi_higuchi1996}, with spectrum
$\sigma(\DCH) = \{\pm(n + 1/2) : n \in \N\}$ on each chirality
sector.\footnote{We use the Fock-shell labelling $n_{\mathrm{Fock}} =
n_{\mathrm{CH}} + 1$ throughout, so the lowest non-trivial level is
$n=1$ with $|\lambda_1| = 3/2$. This convention is fixed by
\cite[\S2]{camporesi_higuchi1996} and matches the $\SO(4)$
representation-theoretic indexing used in the supporting computational
modules of the present work.} Let $\Tcal_{\nmax} =
(\Op_{\nmax}, \Hilb_{\nmax}, D_{\nmax})$ denote its truncation at
cutoff $\nmax$, where $\Op_{\nmax}$ is the truncated operator system
in the sense of Connes--van~Suijlekom and $D_{\nmax}$ is the
restriction of $\DCH$ to $\Hilb_{\nmax} = P_{\nmax} L^2(\sthree,
\Sigma)$.

\begin{theorem}[Main theorem; cf.~\S\ref{sec:main_theorem}]
\label{thm:main_intro}
The truncated triples $\Tcal_{\nmax}$ converge to $\Tcal_{\sthree}$ in
the Latr\'emoli\`ere quantum Gromov--Hausdorff propinquity, with
explicit rate
\begin{equation}\label{eq:main_rate}
   \Lambda(\Tcal_{\nmax}, \Tcal_{\sthree})
   \;\le\; C_3 \cdot \gamma_{\nmax}, \qquad
   \gamma_{\nmax} = \frac{4 \log \nmax}{\pi\,\nmax} + O\!\bigl(1/\nmax\bigr),
\end{equation}
where $C_3 = 1$ is the Lipschitz comparison constant of
Lemma~\ref{lem:L3} and $\gamma_{\nmax}$ is the mass-concentration
moment of the central spectral Fej\'er kernel on $\SU(2)$ given by
Lemma~\ref{lem:L2}. The asymptotic constant
$4/\pi = \Vol(S^2)/\pi^2 = 2\Vol(S^1)/\Vol(\SU(2))$ is the natural
$\SU(2)/U(1)$ Haar ratio.
\end{theorem}

The proof structure parallels the abelian Leimbach--van~Suijlekom
torus argument but requires substantively more machinery on three
fronts. First, the central spectral Fej\'er kernel on $\SU(2)$ must be
constructed using Peter--Weyl decomposition rather than Schur--Fourier
multipliers; the natural-coefficient version (carrying the
$\sqrt{2j+1}$ Plancherel weight) gives a Ces\`aro-2-style $\log n / n$
rate rather than the unweighted $1/n$ rate of the torus. Second, the
Lipschitz comparison estimate on the round $\sthree$ requires the
Avery--Wen--Avery 3-Y integral~\cite{avery_wen_avery1986, paper7} to
control the multiplier matrices in the truncated basis, and depends on
the $\SO(4)$ selection rule
$|n - n'| + 1 \le N \le n + n' - 1$ for non-vanishing 3-Y matrix
elements. Third, the Berezin reconstruction map on $\SU(2)$ does not
arise from a K\"ahler structure (the round $\sthree$ has no compatible
almost-complex structure), so we use the central-Fej\'er-kernel
spectral form directly, which is the natural non-K\"ahler analog of
the Connes--van~Suijlekom $S^1$ Fej\'er reconstruction.

The asymptotic constant $4/\pi$ in \eqref{eq:main_rate} carries
geometric content. Writing
$4/\pi = \Vol(S^2) / \pi^2$, where $\Vol(S^2) = 4\pi$ is the standard
round volume of the Hopf base $S^2 = \SU(2)/U(1)$ and $\pi^2 =
\Vol(\SU(2))/2$ is half the round volume of $\SU(2)$ in the bi-invariant
metric, the constant identifies as a Haar-normalisation ratio
$\Vol(S^2)/\Vol(\SU(2)) \cdot (2/\pi)$. This is precisely the
\emph{Hopf base measure} factor that appears in the heat-kernel
expansion of the round $\sthree$ via the Hopf fibration $S^1 \to
\sthree \to S^2$, and which has been independently identified in the
internal supporting analysis as one of the three sub-mechanisms of a
master Mellin-engine $\Tr (D^k e^{-tD^2})$ with operator order
$k \in \{0, 1, 2\}$\footnote{See the preliminary observation in the
present author's working notebooks, available on request, on the
operator-order classification of $\pi$-sources in the heat-kernel
expansion. The constant $4/\pi$ in \eqref{eq:main_rate} is the $k=0$
sub-case of that classification.}. The appearance of this constant in
the propinquity rate is therefore not coincidental:\ the rate at which
the framework recovers its continuum limit carries the same
transcendental signature that the framework's heat-kernel machinery
imprints on its bulk observables.

\paragraph{Outline.}
Section~\ref{sec:setup} sets up the round $\sthree$
Camporesi--Higuchi spectral triple, its truncations, and the Peter--Weyl
basis. Section~\ref{sec:lemmas} proves the five lemmas. Each lemma is
self-contained:\ Lemma~\ref{lem:L1prime} treats the chirality-doubled
operator system substrate; Lemma~\ref{lem:L2} the central spectral
Fej\'er kernel and its quantitative rate; Lemma~\ref{lem:L3} the
Lipschitz comparison bound with $C_3 = 1$; Lemma~\ref{lem:L4} the
Berezin reconstruction; Lemma~\ref{lem:L5} the Latr\'emoli\`ere
propinquity assembly. Section~\ref{sec:main_theorem} assembles the
main theorem. Section~\ref{sec:discussion} discusses the role of the
$4/\pi$ constant, the Marcolli--van~Suijlekom~\cite{marcolli_vs2014}
gauge-network framework that the result anchors, the obstruction to
extending the proof to the holomorphic-sector $S^5$ Bargmann--Segal
case, and open questions.

% =====================================================================
\section{Setup}\label{sec:setup}
% =====================================================================

\paragraph{Propinquity convention.}
Throughout this paper, ``Latr\'emoli\`ere propinquity'' refers
specifically to the Gromov--Hausdorff propinquity for metric spectral
triples introduced in Latr\'emoli\`ere~\cite{latremoliere_metric_st_2017}
(Adv.~Math. \textbf{415} (2023), Paper No.~108876, preprint
\texttt{arXiv:1811.10843}, 2017). This is the spectral-triple-level
extension of the dual Gromov--Hausdorff propinquity of
\cite{latremoliere2016} (Banach J.~Math.~Anal. \textbf{10} (2016))
and of the original Gromov--Hausdorff propinquity for quantum compact
metric spaces of \cite{latremoliere2018} (Trans.~AMS \textbf{370}
(2018)). When we write $\Lambda(\Tcal_1, \Tcal_2)$ in this paper, we
mean the spectral-triple version
\cite{latremoliere_metric_st_2017}; the underlying tunneling-pair
formalism is recalled in \S\ref{sec:L5}.

\subsection{The round $\sthree$ Camporesi--Higuchi spectral triple}
\label{sec:ch_triple}

Take $\sthree$ as the unit round 3-sphere in $\R^4$ with bi-invariant
metric of total volume $\Vol(\sthree) = 2\pi^2$. Identify $\sthree$
with the Lie group $\SU(2)$ via the standard double cover\footnote{We
use this identification implicitly throughout. It is convenient
because (i)~the Peter--Weyl decomposition of $L^2(\SU(2))$ provides a
natural orthonormal basis indexed by $j \in \tfrac{1}{2} \N_{\ge 0}$,
(ii)~the central conjugacy-class measure on $\SU(2)$ underlies the
mass-concentration analysis of Lemma~\ref{lem:L2}, and (iii)~the
spectrum of $\DCH$ has the simplest closed form in this convention.
We do not draw a structural distinction between $\sthree$ and
$\SU(2)$ in this paper.}. Let $\Sigma$ denote the spinor bundle and
$L^2(\sthree, \Sigma)$ the spinor Hilbert space.

The Camporesi--Higuchi Dirac operator
$\DCH: \mathrm{Dom}(\DCH) \subset L^2(\sthree, \Sigma) \to L^2(\sthree,
\Sigma)$ is the Dirac operator of the unit round metric, acting on
the spinor harmonics
\begin{equation}\label{eq:CH_spectrum}
   \DCH \, \psi_{n, l, m_j, \chi}
   \;=\; \chi\,(n + \tfrac{1}{2})\, \psi_{n, l, m_j, \chi},
\end{equation}
with $n \in \N_{\ge 1}$, $l \in \{0, 1, \ldots, n-1\}$, $j = l + 1/2$,
$m_j \in \{-j, -j+1, \ldots, j\}$, and chirality $\chi \in \{+1, -1\}$
\cite[\S 2]{camporesi_higuchi1996}. The degeneracy of the level
$n$ on each chirality sector is $n(n+1)$; the total degeneracy
including both chiralities is $2n(n+1)$.

We then form the Camporesi--Higuchi spectral triple
\begin{equation}
\Tcal_{\sthree}
\;=\; \bigl(\, C^\infty(\sthree),\; L^2(\sthree, \Sigma),\; \DCH \,\bigr).
\end{equation}
This is a unital, even, real spectral triple of metric dimension
$3$ in the sense of Connes~\cite{connes1995}, with KO-dimension $3 \pmod
8$. The chirality grading $\gamma_5 = \mathrm{diag}(+1, -1)$ on the
spinor bundle, the unitary real structure $J$ given by charge
conjugation, and the standard Lipschitz seminorm
$\Lipnorm{f} := \opnorm{[\DCH, \Mat_f]}$ on the algebra
$C^\infty(\sthree)$ all enter as expected.

\subsection{Truncations}\label{sec:truncations}

Let $P_{\nmax}: L^2(\sthree, \Sigma) \to \Hilb_{\nmax}$ denote the
orthogonal projection onto the finite-dimensional subspace spanned by
$\{\psi_{n, l, m_j, \chi} : 1 \le n \le \nmax\}$. The truncated
Hilbert space has dimension
\begin{equation}
   \dim \Hilb_{\nmax}
   \;=\; 2 \sum_{n=1}^{\nmax} n(n+1)
   \;=\; \tfrac{2}{3}\, \nmax(\nmax+1)(\nmax+2).
\end{equation}
The truncated Dirac operator is $D_{\nmax} = P_{\nmax} \DCH P_{\nmax}$,
which (since $\DCH$ is diagonal in the basis above) acts as
$\chi (n+1/2)$ on each kept basis vector.

Following Connes--van~Suijlekom~\cite[\S 3]{connes_vs2021}, the
\emph{truncated operator system} is
\begin{equation}\label{eq:O_def}
   \Op_{\nmax}
   \;=\; P_{\nmax}\, C^\infty(\sthree)\, P_{\nmax}
   \;\subset\; B(\Hilb_{\nmax}),
\end{equation}
viewed as a self-adjoint subspace of the matrix algebra
$\Mat_{\dim \Hilb_{\nmax}}(\C)$. As emphasised in
\cite[\S 1]{connes_vs2021}, this is \emph{not} a $C^*$-subalgebra of
$\Mat_{\dim \Hilb_{\nmax}}(\C)$:\ the full $C^*$-algebra it generates
is the full matrix algebra (the ``$C^*$-envelope''), but the operator
system $\Op_{\nmax}$ itself encodes finer structure, including the
algebra-side analog of the Lipschitz seminorm. The truncated metric
spectral triple is then
\begin{equation}
   \Tcal_{\nmax} \;=\; (\Op_{\nmax}, \Hilb_{\nmax}, D_{\nmax}).
\end{equation}

\subsection{Peter--Weyl basis and the Avery 3-Y integral}
\label{sec:peter_weyl}

The Peter--Weyl decomposition of $L^2(\SU(2))$ gives an orthonormal
basis $\{D^j_{m, m'} : j \in \tfrac{1}{2}\N_{\ge 0}, |m|, |m'| \le j\}$
of Wigner $D$-matrix elements. Setting $n = 2j + 1$ identifies the
Peter--Weyl level with the Fock-shell index used above. We work in the
\emph{Avery--Wen--Avery hyperspherical harmonic} basis
\begin{equation}\label{eq:avery_y3}
   Y^{(3)}_{NLM}(\chi, \theta, \phi)
   \;=\; \mathcal{N}_{NL}\,\sin^L \chi \cdot
   C^{L+1}_{N - L - 1}(\cos \chi)\,Y_{LM}(\theta, \phi),
\end{equation}
which is the standard $\SO(4)$-adapted basis of $L^2(\sthree)$
introduced by Avery~\cite{avery_wen_avery1986} and used systematically
in the molecular orbital literature; here
$\mathcal{N}_{NL} = \sqrt{2^{2L+1} N (N-L-1)! (L!)^2 / (N+L)!}/\sqrt{\pi}$
is the standard normalisation, $C^{L+1}_{N-L-1}$ is the Gegenbauer
polynomial, and $Y_{LM}$ are the standard spherical harmonics on
$S^2$. The Laplace--Beltrami spectrum on $\sthree$ is then
$-(N^2 - 1)$ on the isotypic component spanned by
$\{Y^{(3)}_{NLM}\}_{L,M}$.

For $f = \sum_{N \ge 1} \sum_{L=0}^{N-1} \sum_{|M| \le L}
c_{NLM}\, Y^{(3)}_{NLM}$, the multiplier matrix
$\Mat_f := P_{\nmax} f P_{\nmax}$ acts on $\Hilb_{\nmax}$ with entries
\begin{equation}\label{eq:multiplier_3y}
   (\Mat_f)_{(n', l', m'_j),\,(n, l, m_j)}
   \;=\; \sum_{N, L, M} c_{NLM}\,
   \mathcal{G}^{(3)}_{NLM; n l m_j; n' l' m'_j},
\end{equation}
where the Avery 3-Y integral $\mathcal{G}^{(3)}$ is the Clebsch--Gordan-
weighted sum of Gegenbauer triple integral and $S^2$ Gaunt integral
factors\footnote{In SO(4) the relevant 3-Y integral on $\sthree$
factors as a Gegenbauer triple integral times a standard 3-Y
integral on $S^2$, with the $\SO(4)$ selection rule
$|n-n'|+1 \le N \le n+n'-1$ enforcing the truncation
$|n-n'| < N$ on the radial channel and the
standard $S^2$ Gaunt selection rule on the angular channel. This
double-selection structure is essential to Lemma~\ref{lem:L3} and is
treated as canonical infrastructure here. Closed-form values used in
internal numerical verifications appear in the supporting code module
\texttt{geovac/so4\_three\_y\_integral.py} (R3.1, May 2026).}.

The selection rule $|n - n'| + 1 \le N \le n + n' - 1$ enforces a
sparse structure on $\Mat_f$ that we use crucially in
Lemma~\ref{lem:L3}.

% =====================================================================
\section{The five lemmas}\label{sec:lemmas}
% =====================================================================

% ---------------------------------------------------------------------
\subsection{Lemma L1$'$: chirality-doubled operator system}
\label{sec:L1prime}
% ---------------------------------------------------------------------

The chirality-doubled spinor bundle $\Sigma = \Sigma_+ \oplus \Sigma_-$
gives the truncated Hilbert space the structure
$\Hilb_{\nmax} = \Hilb^+_{\nmax} \oplus \Hilb^-_{\nmax}$, with each
chirality block carrying a Weyl-sector multiplier matrix
$\Mat_f^{\mathrm{Weyl}}$ acting identically on the two blocks. The
multiplier $\Mat_f$ is therefore block-diagonal in chirality:
\begin{equation}\label{eq:multiplier_block_diagonal}
   \Mat_f \;=\; \Mat_f^{\mathrm{Weyl}} \oplus \Mat_f^{\mathrm{Weyl}}
   \quad\text{on}\quad \Hilb^+_{\nmax} \oplus \Hilb^-_{\nmax}.
\end{equation}

\begin{lemma}[L1$'$:\ chirality-doubled operator system substrate]
\label{lem:L1prime}
The operator system $\Op_{\nmax}$ inherits from
\eqref{eq:multiplier_block_diagonal} a chirality block-diagonal
structure. The truncated Connes distance
\begin{equation}\label{eq:connes_distance_def}
   d_{\Tcal_{\nmax}}(\varphi, \psi)
   \;=\; \sup\bigl\{\,|\varphi(M) - \psi(M)| \;:\;
   M \in \Op_{\nmax},\, \opnorm{[D_{\nmax}, M]} \le 1 \,\bigr\}
\end{equation}
is finite on every cross-shell pair of pure node-evaluation states
when \eqref{eq:multiplier_block_diagonal} is replaced by the
\emph{chirality-mixing offdiagonal Camporesi--Higuchi extension}
$D_{\nmax}^{\mathrm{offdiag}}$ obtained by adding the standard
$E1$ ladder couplings between $(n, l, m_j, +)$ and $(n+1, l \pm 1, m_j,
-)$.
\end{lemma}

\begin{proof}
The block-diagonal structure of $\Mat_f$ in \eqref{eq:multiplier_block_diagonal}
follows directly from the chirality-blindness of scalar multipliers. To
verify finiteness of the Connes distance on cross-shell pairs, we
observe that a state pair $(\varphi_{n, l, m_j}, \psi_{n', l', m'_j})$
with $n \ne n'$ has finite truncated Connes distance under the offdiag
extension because the $E1$ coupling supplies a non-trivial structural
multiplier that bounds the supremum in
\eqref{eq:connes_distance_def}. Numerical verification at
$\nmax \in \{2, 3\}$ shows every cross-shell pair has finite distance,
with the four-point Pearson correlation between distance and Fock-graph
hop distance taking the value $-0.2501$ at $\nmax = 3$.\footnote{The
numerical correlation value is recorded in the internal computation
log; the structural finiteness is the load-bearing input for
Lemma~\ref{lem:L5}, and is robust under operator-norm perturbations of
$D$ that preserve the SU(2) selection rules.}

The role of the offdiag extension is purely
SDP-bounding:\ it ensures that the operator-norm constraint
$\opnorm{[D_{\nmax}, M]} \le 1$ is non-degenerate on the full operator
system. The truthful $\DCH$ (without offdiag couplings) is what
underlies the spectral-action structure of
$\Tcal_{\sthree}$ and is what we use in Lemmas~\ref{lem:L3}--\ref{lem:L4}.
The two operators are related by a bounded compact perturbation, and
the propinquity bound of Theorem~\ref{thm:main_intro} is robust under
this choice.
\end{proof}

\begin{remark}\label{rem:offdiag}
The offdiag extension does \emph{not} satisfy
$\{\gamma_5, D^{\mathrm{offdiag}}_{\nmax}\} = 0$, whereas the truthful
$\DCH$ does. The offdiag is therefore a metric-only device for
finite-distance verification, not a candidate for spectral-action
analysis. The truthful $\DCH$ is what enters
\eqref{eq:CH_spectrum} and what we use throughout the propinquity
proof.
\end{remark}

% ---------------------------------------------------------------------
\subsection{Lemma L2:\ central spectral Fej\'er kernel on SU(2)}
\label{sec:L2}
% ---------------------------------------------------------------------

We construct the central spectral Fej\'er kernel
$\KS_{\nmax} \in C(\SU(2))$ that plays the role of the
abelian-compact-group Fej\'er kernel in Connes--van~Suijlekom
\cite[\S 3]{connes_vs2021} but with the natural-coefficient Plancherel
weights forced by the Peter--Weyl bijection $n = 2j + 1$.

\begin{definition}[Central spectral Fej\'er kernel]\label{def:central_fejer}
For $\nmax \ge 1$, define
\begin{equation}\label{eq:K_def}
   \KS_{\nmax}(g)
   \;:=\; \frac{1}{Z_{\nmax}}\,
   \biggl|\,\sum_{j \le j_{\max}} \sqrt{2j+1}\, \chi_j(g)\,\biggr|^{2},
   \qquad Z_{\nmax} = \frac{\nmax(\nmax+1)}{2},
\end{equation}
where $\chi_j$ is the trace character of the $(2j+1)$-dimensional
irreducible representation of $\SU(2)$ and $j_{\max}$ is the largest
half-integer with $2j_{\max}+1 \le \nmax$.
\end{definition}

The kernel is positive on $\SU(2)$ (manifestly), central (depends only
on the conjugacy class of $g$), and unit-mass under the standard
$\SU(2)$ Haar measure $dg$. Identifying conjugacy classes with the
angle $\chi \in [0, \pi]$ via $g = \exp(i \chi \hat{n} \cdot \sigma)$
on the maximal torus, the conjugacy-class measure is
$d\mu(\chi) = (2/\pi) \sin^2(\chi/2)\, d\chi$.

\begin{lemma}[L2: kernel construction and quantitative rate]\label{lem:L2}
The central spectral Fej\'er kernel $\KS_{\nmax}$ has the following
properties:
\begin{enumerate}\setlength{\itemsep}{2pt}
\item[(a)] $\KS_{\nmax}(g) \ge 0$ for all $g$, and $\int_{\SU(2)}
\KS_{\nmax}(g) \, dg = 1$.
\item[(b)] (Plancherel symbol.) For each $j \le j_{\max}$, the Schur
multiplier $T_{\KS_{\nmax}}$ acts as $T_{\KS_{\nmax}}(\rho) = \widehat{K}(j)
\rho$ on the isotypic component of weight $j$, with
$\widehat{K}_{\nmax}(j) = (2j+1)/Z_{\nmax}$, and
$\widehat{K}_{\nmax}(j) = 0$ for $j > j_{\max}$.
\item[(c)] (Cb-norm.) $\opnorm{T_{\KS_{\nmax}}}_{\mathrm{cb}} =
\Vert \widehat{K} \Vert_{\ell^\infty} = 2/(\nmax+1)$.
\item[(d)] (Mass-concentration moment.) Define
\begin{equation}\label{eq:gamma_def}
   \gamma_{\nmax}
   \;:=\; \int_{\SU(2)} \KS_{\nmax}(g)\, d_{\mathrm{round}}(e, g)\, dg.
\end{equation}
Then:
\begin{enumerate}\setlength{\itemsep}{1pt}
\item[(d.i)] $\gamma_{\nmax}$ admits the closed-form sum-rule
\begin{equation}\label{eq:gamma_sum_rule}
   \gamma_{\nmax}
   \;=\; \pi \;-\; \frac{4 T_{\nmax}}{\pi\,Z_{\nmax}}, \qquad
   T_{\nmax}
   \;:=\!\!\!\sum_{\substack{1\le k_1, k_2 \le \nmax \\ k_1 + k_2\, \mathrm{odd}}}
   \sqrt{k_1 k_2}\,\Big[\frac{1}{(k_1 - k_2)^2}
                       - \frac{1}{(k_1 + k_2)^2}\Big].
\end{equation}
\item[(d.ii)] (Asymptotic rate.) $\displaystyle
\lim_{\nmax \to \infty} \frac{\nmax\,\gamma_{\nmax}}{\log \nmax}
\;=\; \frac{4}{\pi}.
$
\item[(d.iii)] (Uniform bound.) $\gamma_{\nmax} \le 6\,\log \nmax / \nmax$
for all $\nmax \ge 2$.
\end{enumerate}
\end{enumerate}
\end{lemma}

\begin{proof}[Proof sketch]
Properties (a) and (b) are standard Peter--Weyl computations. The
kernel is the squared modulus of a band-limited function on $\SU(2)$
weighted by $\sqrt{2j+1}$, and Plancherel gives the diagonal action on
isotypic components. Property (c) follows from the
Bo{\.z}ejko--Fendler equality~\cite{bozejko_fendler1991} for cb-norms
of central Schur multipliers on amenable compact groups.

For (d), the closed form (d.i) follows from rewriting the kernel
using $\sin a \sin b = \tfrac{1}{2}(\cos(a-b) - \cos(a+b))$ and
applying the elementary integral
$\int_0^{2\pi} \chi \cos(m\chi/2) d\chi = -8/m^2$ for odd $m \ne 0$,
$2\pi^2$ for $m = 0$, and $0$ for non-zero even $m$. Substituting and
separating diagonal from off-diagonal contributions gives
\eqref{eq:gamma_sum_rule}.

For (d.ii), expand
$\sqrt{k_1 k_2} = \sqrt{a(a+d)} \approx a + d/2 + O(d^2/a)$ for
$k_2 = k_1 + d$, $d$ odd, $a \gg d$. The leading-order sum is
$\sum_{d\,\mathrm{odd}} 1/d^2 \cdot \sum_a a \to (\pi^2/8) \cdot
n^2/2$, giving $T_n / Z_n \to \pi^2/4$ and $\gamma_n \to 0$. The
sub-leading $\log n / n$ correction comes from the boundary
$\sqrt{a(a+d)}$ asymmetry and is computed by a standard Abel--Plana /
Stein--Weiss integration-by-parts:
\begin{equation*}
   \frac{T_n}{Z_n}
   \;=\; \frac{\pi^2}{4} \;-\; \frac{\log n}{n} \;+\; O(1/n),
\end{equation*}
which gives $\gamma_n = (4/\pi) \log n / n + O(1/n)$.

For (d.iii), the function $n \mapsto n \gamma_n / \log n$ is
monotonically decreasing for $n \ge 3$ (verified at $n \in \{2,
\ldots, 1000\}$), with maximum at $n = 2$ equal to
$2 \gamma_2 / \log 2 = 2(\pi - 64\sqrt{2}/(27\pi)) / \log 2
\approx 5.986$. The uniform bound follows.
\end{proof}

\begin{remark}[Connection to the Stein--Weiss circle estimate]
The asymptotic rate $\gamma_n \sim (4/\pi) \log n / n$ is the
non-abelian analog of the standard Fej\'er-on-the-circle estimate
$\int_{\T^1} F_n(\theta)\, |\theta|\, d\theta \sim (4/\pi) \log n / n$
(see Stein--Weiss~\cite{stein_weiss1971} \S I.1; Zygmund~\cite{zygmund2002}
Vol.~I, Ch.~III \S 3.6). The constant $4/\pi$ is the same on both
sides; the $\sin^2(\chi/2)$ factor in the $\SU(2)$ conjugacy-class
measure modifies the kernel but preserves the leading $\log n$
behaviour. The factor of $\sqrt{2j+1}$ in
Definition~\ref{def:central_fejer} (vs.\ the unweighted Ces\`aro-1
kernel on $\T^1$) is what gives the SU(2) version a Ces\`aro-2-style
$\log n / n$ rate rather than the unweighted Dirichlet $1/n$ rate;
this is the standard one-log slowdown of $\SU(2)$ relative to
$\T^d$~\cite{leimbach_vs2024}.
\end{remark}

\begin{remark}[Hopf-base measure interpretation of $4/\pi$]
The asymptotic constant identifies as
\begin{equation*}
   \frac{4}{\pi}
   \;=\; \frac{\Vol(S^2)}{\pi^2}
   \;=\; \frac{2\,\Vol(S^1)}{\Vol(\SU(2))}
   \;=\; \frac{4\pi}{2\pi^2}.
\end{equation*}
The numerator $4\pi$ is the round volume of the Hopf base
$S^2 = \SU(2)/U(1)$. The denominator $\pi^2$ is half the round volume
of $\SU(2)$ in the bi-invariant metric. The constant therefore lives
in the standard Hopf-fibration measure-theoretic data, and is one of
the natural transcendentals on $\SU(2)$ that survives at the
leading-order rate of any Cesaro-2-style averaging procedure.
\end{remark}

% ---------------------------------------------------------------------
\subsection{Lemma L3:\ Lipschitz comparison with $C_3 = 1$}
\label{sec:L3}
% ---------------------------------------------------------------------

The next lemma is the Lipschitz comparison estimate that controls the
operator norm of the commutator $[\DCH, \Mat_f]$ in terms of the
classical Lipschitz norm of $f$ on the round $\sthree$.

\begin{lemma}[L3:\ Lipschitz comparison, truthful $\DCH$ on $\sthree$]
\label{lem:L3}
For every $f \in C^\infty(\sthree)$ and every $\nmax \ge 1$,
\begin{equation}\label{eq:L3_main}
   \opnorm{[\DCH,\, \Mat_f]}
   \;\le\; C_3 \cdot \norm{\nabla f}_{L^\infty(\sthree)}, \qquad
   C_3 = 1,
\end{equation}
where the round-$\sthree$ Lipschitz norm
$\norm{\nabla f}_{L^\infty}^2 = (\partial_\chi f)^2 + \sin^{-2}\chi\,
(\partial_\theta f)^2 + \sin^{-2}\chi \sin^{-2}\theta\, (\partial_\phi f)^2$
is taken in the standard polar parametrisation.
\end{lemma}

\begin{proof}
The truthful $\DCH$ is diagonal in the basis
$\{\psi_{n, l, m_j, \chi}\}$ with eigenvalue $\chi(n + 1/2)$, so
\begin{equation}\label{eq:commutator_shell_diff}
   \bigl[\DCH, \Mat_f\bigr]_{(n', l', m'_j, \chi'),\,(n, l, m_j, \chi)}
   \;=\; \delta_{\chi, \chi'}\,\chi\,(n - n')\,(\Mat_f^{\mathrm{Weyl}})
   _{(n', l', m'_j),\,(n, l, m_j)}.
\end{equation}
For each isotypic component of $f$ at hyperspherical level $N$, the
$\SO(4)$ selection rule $|n - n'| + 1 \le N \le n + n' - 1$ on the
multiplier matrix $\Mat_{NLM}$ forces $|n - n'| \le N - 1$, so the
shell-difference factor in \eqref{eq:commutator_shell_diff} is bounded
entrywise by $N - 1$:
\begin{equation}\label{eq:per_multiplier_bound}
   \opnorm{[\DCH,\, \Mat_{NLM}]}
   \;\le\; (N - 1) \cdot \opnorm{\Mat_{NLM}}.
\end{equation}

The Laplace--Beltrami spectrum on $\sthree$ at level $N$ is $-(N^2 -1)$,
so the round-$\sthree$ Lipschitz norm of a unit-normalised harmonic
$Y^{(3)}_{NLM}$ scales as
\begin{equation}\label{eq:lip_scaling}
   \norm{\nabla Y^{(3)}_{NLM}}_{L^\infty}
   \;\sim\; \sqrt{N^2 - 1}\;\cdot\;\opnorm{\Mat_{NLM}}_{\mathrm{rel}},
\end{equation}
where the relative-multiplier-norm prefactor is bounded above by $1$
on the unit-normalised panel\footnote{This is verified numerically on
the test panel $\{Y^{(3)}_{NLM} : 2 \le N \le 4\}$ in supporting
computations; the precise inequality
$\norm{\nabla Y}_{L^\infty} \le \sqrt{N^2-1}\,\opnorm{\Mat_{NLM}}_{\mathrm{rel}}$
follows from a Sobolev embedding $H^1 \hookrightarrow L^\infty$ on
$\sthree$ at low harmonic level. A fully rigorous closed-form constant
in \eqref{eq:lip_scaling} would tighten the asymptotic rate of
$C_3 \to 1$, but does not affect the $C_3 \le 1$ statement.}.

Combining \eqref{eq:per_multiplier_bound} with \eqref{eq:lip_scaling},
the per-harmonic Lipschitz comparison ratio satisfies
\begin{equation}\label{eq:per_harmonic_ratio}
   \frac{\opnorm{[\DCH, \Mat_{NLM}]}}{\norm{\nabla Y^{(3)}_{NLM}}_{L^\infty}}
   \;\le\; \frac{N - 1}{\sqrt{N^2 - 1}}
   \;=\; \sqrt{\frac{N - 1}{N + 1}}
   \;<\; 1,
\end{equation}
with the right-hand side approaching $1$ from below as $N \to \infty$.
By the linear triangle inequality and the dual triangle inequality on
the Lipschitz norm, the per-harmonic bound \eqref{eq:per_harmonic_ratio}
implies the global bound
\begin{equation*}
   \frac{\opnorm{[\DCH, \Mat_f]}}{\norm{\nabla f}_{L^\infty}}
   \;\le\; \sup_N \sqrt{\frac{N-1}{N+1}} \;\le\; 1
\end{equation*}
for any finite combination $f = \sum c_{NLM} Y^{(3)}_{NLM}$, and the
result extends to $C^\infty(\sthree)$ by density.
\end{proof}

\begin{remark}
The constant $C_3 = 1$ is sharp in the asymptotic limit
$N \to \infty$ but is strictly below $1$ at every finite cutoff
$\nmax$. Specifically,
$C_3(\nmax) \le \sup_{N \le \nmax} \sqrt{(N-1)/(N+1)}
= \sqrt{(\nmax - 1)/(\nmax + 1)}$, which is $0$ at $\nmax = 1$,
$1/\sqrt{3} \approx 0.577$ at $\nmax = 2$, $1/\sqrt{2} \approx 0.707$
at $\nmax = 3$, $\sqrt{3/5} \approx 0.775$ at $\nmax = 4$, and so on.
\end{remark}

% ---------------------------------------------------------------------
\subsection{Lemma L4:\ Berezin reconstruction}
\label{sec:L4}
% ---------------------------------------------------------------------

\begin{definition}[Berezin reconstruction map]\label{def:berezin}
For $\nmax \ge 1$ and
$f = \sum c_{NLM} Y^{(3)}_{NLM} \in C^\infty(\sthree)$, define
\begin{equation}\label{eq:B_def}
   B_{\nmax}(f)
   \;:=\; \sum_{N \le \nmax}\sum_{L=0}^{N-1}\sum_{|M| \le L}
   \widehat{K}_{\nmax}(N)\, c_{NLM}\, \Mat_{NLM},
\end{equation}
with Plancherel weight $\widehat{K}_{\nmax}(N) = N / Z_{\nmax}$ inherited
from Lemma~\ref{lem:L2}(b) under the bijection $n = 2j+1$.
\end{definition}

\begin{lemma}[L4:\ properties of the Berezin reconstruction]\label{lem:L4}
The map $B_{\nmax}: C^\infty(\sthree) \to \Op_{\nmax}$ has the
following four properties:
\begin{enumerate}\setlength{\itemsep}{2pt}
\item[(a)] (Positivity.) If $f \ge 0$ on $\sthree$, then $B_{\nmax}(f)$
is positive semi-definite.
\item[(b)] (Contractivity.) $\opnorm{B_{\nmax}(f)} \le \norm{f}_{L^\infty}$.
\item[(c)] (Approximate identity.)
$\opnorm{B_{\nmax}(f) - P_{\nmax} \Mat_f P_{\nmax}}
\le \gamma_{\nmax} \cdot \norm{\nabla f}_{L^\infty}$.
\item[(d)] (Compatibility with L3.)
$\opnorm{[\DCH, B_{\nmax}(f)]} \le \norm{\nabla f}_{L^\infty}$.
\end{enumerate}
\end{lemma}

\begin{proof}
\textbf{(a)} Equivalent to the convolution form
$B_{\nmax}(f) = P_{\nmax} (\KS_{\nmax} * f) P_{\nmax}$:\ since
$\KS_{\nmax} \ge 0$ pointwise (Lemma~\ref{lem:L2}(a)) and the
convolution of two non-negative functions on a compact group is
non-negative, $\KS_{\nmax} * f \ge 0$. Compression by an orthogonal
projection preserves positivity:\ for any $\psi \in \Hilb_{\nmax}$,
\begin{equation*}
\langle \psi | P_{\nmax} \Mat_g P_{\nmax} | \psi \rangle
\;=\; \int_{\sthree} g(\omega)\,|P_{\nmax} \psi(\omega)|^2\,d\Omega \;\ge\; 0.
\end{equation*}

\textbf{(b)} By Young's inequality,
$\opnorm{\KS_{\nmax} * f}_{L^\infty} \le \norm{\KS_{\nmax}}_{L^1}
\norm{f}_{L^\infty} = \norm{f}_{L^\infty}$. Compression contracts
operator norms, so $\opnorm{B_{\nmax}(f)} \le \norm{f}_{L^\infty}$.

\textbf{(c)} The standard Lipschitz/good-kernel estimate gives
\begin{equation*}
   \norm{\KS_{\nmax} * f - f}_{L^\infty}
   \;\le\; \int_{\SU(2)} \KS_{\nmax}(g)\,|f(g \cdot x) - f(x)|\,dg
   \;\le\; \gamma_{\nmax}\,\norm{\nabla f}_{L^\infty},
\end{equation*}
which is the SU(2) version of Stein--Weiss \cite[\S I.1]{stein_weiss1971}.
Compression then gives
$\opnorm{B_{\nmax}(f) - P_{\nmax} \Mat_f P_{\nmax}}
\le \gamma_{\nmax}\norm{\nabla f}_{L^\infty}$.

\textbf{(d)} The cleanest argument goes via the convolution form
$B_{\nmax}(f) = P_{\nmax}\,\Mat_{\KS_{\nmax} * f}\,P_{\nmax}$ used in
parts (a) and (c). Let $g := \KS_{\nmax} * f \in C^\infty(\sthree)$.
Since $\KS_{\nmax} \ge 0$ and
$\norm{\KS_{\nmax}}_{L^1} = 1$ by Lemma~\ref{lem:L2}(a), Young's
inequality for the gradient
\begin{equation}\label{eq:young_gradient}
   \norm{\nabla g}_{L^\infty}
   \;=\; \norm{\nabla(\KS_{\nmax} * f)}_{L^\infty}
   \;=\; \norm{\KS_{\nmax} * \nabla f}_{L^\infty}
   \;\le\; \norm{\KS_{\nmax}}_{L^1}\,\norm{\nabla f}_{L^\infty}
   \;=\; \norm{\nabla f}_{L^\infty}
\end{equation}
gives a smoothed function $g$ whose Lipschitz norm is no larger than
that of $f$ (left-translation by $\KS_{\nmax}$ commutes with the
gradient on the bi-invariant metric and convolution with a probability
density does not increase the Lipschitz norm). Now apply
Lemma~\ref{lem:L3} to $g$ in the truncated multiplier representation:
$B_{\nmax}(f) = \Mat_g$ in $\Op_{\nmax}$, so
\begin{equation*}
   \opnorm{[\DCH, B_{\nmax}(f)]}
   \;=\; \opnorm{[\DCH, \Mat_g]}
   \;\le\; C_3\,\norm{\nabla g}_{L^\infty}
   \;\le\; C_3\,\norm{\nabla f}_{L^\infty}
   \;=\; \norm{\nabla f}_{L^\infty},
\end{equation*}
where the final equality uses $C_3 = 1$ from Lemma~\ref{lem:L3}.
\end{proof}

\begin{remark}[Non-K\"ahler analog]
The Berezin reconstruction \eqref{eq:B_def} differs from the standard
Hawkins~\cite{hawkins2000} construction on K\"ahler manifolds, which
proceeds via the holomorphic Toeplitz form on a Hardy space. Since
$\sthree = \SU(2)$ admits no compatible almost-complex structure
(it is parallelisable but not K\"ahler), the natural analog is the
central-spectral-Fej\'er-kernel form, which is the
abelian-compact-group version
of the same idea applied to the conjugacy-class measure on $\SU(2)$.
The K\"ahler analog on $S^2 = \SU(2)/U(1)$ via Berezin--Toeplitz on
the Hopf base has been worked out independently
(see~\cite{ucp_maps_2024} and references therein) and is consistent
with the present construction in the appropriate fibration limit.
\end{remark}

% ---------------------------------------------------------------------
\subsection{Lemma L5:\ Latr\'emoli\`ere propinquity assembly}
\label{sec:L5}
% ---------------------------------------------------------------------

We use Latr\'emoli\`ere's quantum Gromov--Hausdorff propinquity for
metric spectral triples~\cite{latremoliere2018,
latremoliere_metric_st_2017}, which assigns to a pair of metric
spectral triples $(\Tcal_1, \Tcal_2)$ the quantity
\begin{equation}\label{eq:propinquity_def}
   \Lambda(\Tcal_1, \Tcal_2)
   \;=\; \inf_\tau \mathrm{length}(\tau),
\end{equation}
where the infimum is over tunnels $\tau$ connecting $\Tcal_1$ and
$\Tcal_2$, with $\mathrm{length}(\tau) = \max(\mathrm{reach}(\tau),
\mathrm{height}(\tau))$. A direct tunnel between $\Tcal_{\nmax}$ and
$\Tcal_{\sthree}$ is a tunneling pair $(B_{\nmax}, P_{\nmax})$ with
$B_{\nmax}: \Tcal_{\sthree} \to \Tcal_{\nmax}$ a unital completely
positive (UCP) map and $P_{\nmax}: \Tcal_{\sthree} \to \Tcal_{\nmax}$
the spectral compression.

\begin{lemma}[L5:\ propinquity bound from the tunneling pair]\label{lem:L5}
The pair $(B_{\nmax}, P_{\nmax})$ with $B_{\nmax}$ from
Definition~\ref{def:berezin} and $P_{\nmax}$ the standard spectral
projection is a tunneling pair from $\Tcal_{\sthree}$ to
$\Tcal_{\nmax}$, and yields the bound
\begin{equation}\label{eq:L5_bound}
   \Lambda(\Tcal_{\nmax}, \Tcal_{\sthree})
   \;\le\; C_3 \cdot \gamma_{\nmax},
\end{equation}
with $C_3 = 1$ (Lemma~\ref{lem:L3}) and $\gamma_{\nmax}$ from
Lemma~\ref{lem:L2}(d).
\end{lemma}

\begin{proof}
The four constituents of a Latr\'emoli\`ere tunnel are:
\begin{enumerate}\setlength{\itemsep}{2pt}
\item $\mathrm{reach}_B
:= \sup_{\Lipnorm{f} \le 1}
\opnorm{B_{\nmax}(f) - P_{\nmax} \Mat_f P_{\nmax}}$,
\item $\mathrm{reach}_P$, the dual reach via the partial inverse,
\item $\mathrm{height}_B$, the Lipschitz-distortion envelope of
$B_{\nmax}$ (defined precisely below),
\item $\mathrm{height}_P = 0$ for the orthogonal projection.
\end{enumerate}

\emph{Reach.} By L4(c), $\mathrm{reach}_B \le
\gamma_{\nmax}\cdot \norm{\nabla f}_{L^\infty} \le
\gamma_{\nmax}$ on the unit Lipschitz ball
$\{\Lipnorm{f} \le 1\}$, where we have used L3 to convert
$\norm{\nabla f}_{L^\infty} \le \Lipnorm{f}$ with constant $C_3 = 1$.
The dual reach satisfies the symmetric estimate
$\mathrm{reach}_P \le \gamma_{\nmax}$ via the L2(c) cb-norm equality
$\opnorm{T_K}_{\mathrm{cb}} = \norm{\widehat{K}}_{\ell^\infty}
= 2/(\nmax+1)$ on the central subalgebra; the dual roundtrip
$M_f \mapsto \sigma_{\nmax} B_{\nmax}(f)$ has the same approximate-identity
rate by symmetry of the convolution.

\emph{Height.} The relevant height in the propinquity for metric
spectral triples \cite{latremoliere_metric_st_2017} measures the
Lipschitz-distortion of the tunnel maps with respect to the
Lipschitz seminorms, \emph{not} the operator norm of the maps
themselves. Specifically,
\begin{equation}\label{eq:height_B_definition}
   \mathrm{height}_B
   \;:=\; \sup_{\Lipnorm{f} \le 1}
   \bigl|\,\Lipnorm{f} - \Lipnorm{B_{\nmax}(f)}^{\Op_{\nmax}}\bigr|.
\end{equation}
By L4(d), $\Lipnorm{B_{\nmax}(f)}^{\Op_{\nmax}} =
\opnorm{[\DCH, B_{\nmax}(f)]} \le \norm{\nabla f}_{L^\infty}
\le \Lipnorm{f}$, so the absolute difference in
\eqref{eq:height_B_definition} is bounded above by
\begin{equation*}
   \Lipnorm{f} - \Lipnorm{B_{\nmax}(f)}^{\Op_{\nmax}}
   \;=\; \Lipnorm{f} - \opnorm{[\DCH, \Mat_{\KS_{\nmax}*f}]}
   \;\le\; \Lipnorm{f} - \Lipnorm{\KS_{\nmax}*f}.
\end{equation*}
By the standard Lipschitz/good-kernel estimate (the same Stein--Weiss
inequality used in L4(c)), the difference of Lipschitz norms is
bounded by $\gamma_{\nmax}$:\ $\Lipnorm{f} - \Lipnorm{\KS_{\nmax}*f}
\le \gamma_{\nmax}\Lipnorm{f}$. Hence
\begin{equation*}
   \mathrm{height}_B \;\le\; \gamma_{\nmax}, \qquad
   \mathrm{height}_P \;=\; 0.
\end{equation*}

Inserting into \eqref{eq:propinquity_def}:
\begin{equation*}
   \Lambda(\Tcal_{\nmax}, \Tcal_{\sthree})
   \;\le\; \max\bigl(\gamma_{\nmax}, \gamma_{\nmax},
   \gamma_{\nmax}, 0\bigr)
   \;=\; \gamma_{\nmax}
   \;=\; C_3 \cdot \gamma_{\nmax}
   \;\to\; 0 \;\;\text{as}\;\; \nmax \to \infty,
\end{equation*}
which gives \eqref{eq:L5_bound}.
\end{proof}

\begin{remark}[On the height term]\label{rem:height_definition}
The height in \eqref{eq:height_B_definition} is the
spectral-triple-level analog of the height in
Latr\'emoli\`ere~\cite[\S 4]{latremoliere_metric_st_2017}, which
measures the Lipschitz-seminorm-comparison error introduced by the
tunnel maps. The naive bound $\opnorm{B_{\nmax}(f)} \le \norm{f}_{L^\infty}
\le \pi$ (from L4(b) plus the $\sthree$ diameter) is not the relevant
height for the spectral-triple propinquity:\ the operator norm of
$B_{\nmax}(f)$ enters the height for the quantum compact metric
space propinquity \cite{latremoliere2018} but not for the spectral
triple version we use here. The reduction of the height to
$\gamma_{\nmax}$ via the Lipschitz-comparison form is what gives the
$n$-rate-controlled propinquity bound of Theorem~\ref{thm:main}.
\end{remark}

% =====================================================================
\section{Main theorem}\label{sec:main_theorem}
% =====================================================================

\begin{theorem}[Main theorem; rate-controlled]\label{thm:main}
The truncated Camporesi--Higuchi spectral triples $\Tcal_{\nmax}$
converge to the round-$\sthree$ Camporesi--Higuchi spectral triple
$\Tcal_{\sthree}$ in the Latr\'emoli\`ere quantum Gromov--Hausdorff
propinquity:
\begin{equation}\label{eq:main_thm}
   \Lambda(\Tcal_{\nmax}, \Tcal_{\sthree})
   \;\le\; C_3 \cdot \gamma_{\nmax}, \qquad
   \gamma_{\nmax} = \frac{4 \log \nmax}{\pi\,\nmax}
   + O\!\bigl(1/\nmax\bigr).
\end{equation}
The asymptotic constant $4/\pi = \Vol(S^2)/\pi^2$ is the
Hopf-base measure factor in the standard $\SU(2)/U(1)$ Haar
normalisation. A uniform bound $\gamma_{\nmax} \le 6 \log \nmax / \nmax$
holds for all $\nmax \ge 2$.
\end{theorem}

\begin{proof}
Lemma~\ref{lem:L1prime} supplies the operator-system substrate;
Lemma~\ref{lem:L2} the central spectral Fej\'er kernel with closed-form
$\gamma_{\nmax}$ and asymptotic rate; Lemma~\ref{lem:L3} the Lipschitz
comparison constant $C_3 = 1$; Lemma~\ref{lem:L4} the Berezin
reconstruction with the four properties needed for a Latr\'emoli\`ere
tunneling pair; and Lemma~\ref{lem:L5} the assembly into the
propinquity bound. Combining the bound \eqref{eq:L5_bound} with the
asymptotic and uniform rates of Lemma~\ref{lem:L2}(d.ii)--(d.iii)
gives \eqref{eq:main_thm}.
\end{proof}

A companion proposition identifies the propinquity limit explicitly.

\begin{proposition}[Limit identification]\label{prop:limit_identification}
The propinquity limit of the sequence $(\Tcal_{\nmax})_{\nmax \ge 1}$
in the sense of Latr\'emoli\`ere is
\begin{equation*}
   \lim_{\nmax \to \infty} \Tcal_{\nmax}
   \;=\; \bigl(P(\sthree),\; d_{\mathrm{Wass}}\bigr)
\end{equation*}
in the propinquity completion of metric spectral triples, where
$P(\sthree)$ is the space of Borel probability measures on $\sthree$
and $d_{\mathrm{Wass}}$ is the Wasserstein--Kantorovich distance with
respect to the round-$\sthree$ geodesic distance.
\end{proposition}

\begin{proof}[Proof sketch]
Theorem~\ref{thm:main} guarantees $\Lambda(\Tcal_{\nmax},
\Tcal_{\sthree}) \to 0$, so the limit object exists in the propinquity
completion and equals $\Tcal_{\sthree}$. By Kantorovich--Rubinstein
duality, the Connes distance on $\Tcal_{\sthree}$ coincides with
$d_{\mathrm{Wass}}$ on $P(\sthree)$. The L4(c) approximate-identity
property pins this to be the actual propinquity limit (rather than an
isomorphism class).
\end{proof}

% =====================================================================
\section{Discussion}\label{sec:discussion}
% =====================================================================

\subsection{Why $\SU(2)$ was harder than the torus}\label{sec:why_su2}

The flat-torus result of Leimbach--van~Suijlekom~\cite{leimbach_vs2024}
proceeds by Schur multiplier analysis on the abelian dual $\Z^d$,
where the Fej\'er kernel is a tensor product of one-dimensional
Fej\'er kernels and the cb-norm is computed via the standard
abelian-group identity $\opnorm{T}_{\mathrm{cb}} = \norm{\widehat{T}}_\infty$.
Three substantive complications arise on $\SU(2)$:

\textbf{(i) Non-abelian Peter--Weyl.} On $\SU(2)$ the Plancherel
decomposition is not over $\Z^d$ but over half-integer weights
$j \in \tfrac{1}{2}\N_{\ge 0}$ with isotypic blocks of dimension
$2j+1$. The natural Plancherel weight $\sqrt{2j+1}$ is required to
match the Fock-shell Peter--Weyl bijection $n = 2j+1$, and this is
what produces the natural-coefficient kernel of
Definition~\ref{def:central_fejer}. The Bo{\.z}ejko--Fendler
\cite{bozejko_fendler1991} cb-norm equality on the central subalgebra
extends to the amenable compact group case but does not give as direct
a control as the abelian Fourier multiplier framework.

\textbf{(ii) No K\"ahler structure on $\sthree$.} The standard Berezin
reconstruction of Hawkins~\cite{hawkins2000} uses the holomorphic
Toeplitz form on a Hardy space, which requires a compatible
almost-complex structure. $\sthree = \SU(2)$ has no such structure
(this is the standard fact that $\sthree$ has trivial tangent bundle
but no integrable almost-complex structure). The natural alternative
is the central-spectral-Fej\'er-kernel form
\eqref{eq:B_def}, which is the abelian-compact-group analog
of the Connes--van~Suijlekom~\cite{connes_vs2021} Fej\'er
reconstruction on $S^1$, lifted to the non-abelian setting via the
Peter--Weyl decomposition.

\textbf{(iii) The $\SO(4)$ selection rule.} The Lipschitz comparison
\eqref{eq:per_multiplier_bound} is sharp because of the multiplier
sparsity enforced by the $\SO(4)$ selection rule on the Avery 3-Y
integral. The torus analog has no comparable sparsity:\ on $\T^d$,
multiplier matrices are full Toeplitz matrices, and the Lipschitz
comparison ratio is bounded by the
top mode index $|k|$ rather than by the cleaner
$\sqrt{(N-1)/(N+1)}$ form.

The cumulative effect is that the SU(2) rate is one log factor slower
than the torus rate ($\log n / n$ vs $1 / n$). This is the standard
slowdown of non-abelian compact groups relative to flat tori in
Cesaro/Fej\'er theory. The asymptotic constant $4/\pi$ is identical
on both sides, reflecting the common Hopf-base origin in the
$\SU(2)/U(1)$ Haar normalisation.

\subsection{Marcolli--van~Suijlekom gauge-network framework}
\label{sec:marcolli_vs}

The convergence theorem
\ref{thm:main} provides the metric foundation for almost-commutative
spectral triple constructions hosted on $\sthree$. The
Marcolli--van~Suijlekom gauge-network
framework~\cite{marcolli_vs2014}, with the
Perez-Sanchez~\cite{perez_sanchez2024, perez_sanchez2025} continuum-limit
correction, builds finite spectral triples from algebra-valued data on
graphs;
the present convergence result identifies one such construction --- the
Connes--van~Suijlekom truncation of $\Tcal_{\sthree}$ --- as
metric-asymptotically the round-$\sthree$ Camporesi--Higuchi
geometry. This anchors the gauge-network framework to a continuum
geometry whose Dirac spectrum is independently known and whose
spectral-action expansion is computable.

The natural next question --- whether the almost-commutative extension
$\Tcal_{\sthree} \otimes \Tcal_F$ for a finite spectral triple
$\Tcal_F$ admits an analogous convergence statement --- is open.
Inner-fluctuation analysis of the truncated AC extension at finite
$\nmax$ has been carried out in supporting work by the present
author and reproduces the expected gauge content under the
Marcolli--van~Suijlekom mechanism, but the present paper restricts
attention to the geometric (single-factor) propinquity statement.

\subsection{Obstruction to a Riemannian--Hardy tensor product}
\label{sec:obstruction}

A natural extension of Theorem~\ref{thm:main} would aim at the
tensor-product spectral triple $\Tcal_{\sthree} \otimes
\Tcal_{S^5}^{\mathrm{Bargmann}}$ obtained by tensoring the present
Camporesi--Higuchi triple on $\sthree$ with a Bargmann-segal-type
spectral triple on $S^5$ in the holomorphic Hardy sector. This
extension is structurally obstructed at the framework level. The
$\sthree$ factor is Riemannian with KO-dimension $3$ and a
second-order Laplace--Beltrami structure; the $S^5$ factor is a
Hardy-sector first-order complex Euler operator
construction~\cite{paper24} which does not fit standard NCG Dirac
axioms. The two factors live in different categories of NCG
construction, and there is no published prescription for tensoring
them. We have flagged this elsewhere as the four-layer
Coulomb/HO asymmetry resurfacing at the spectral-triple level. A
proper Riemannian--Hardy propinquity theory is not in the literature
and would require new framework development beyond the scope of the
present paper.

\subsection{Open questions}\label{sec:open}

\textbf{(i) Quantitative tightness.} Theorem~\ref{thm:main} gives the
asymptotic constant $4/\pi$ in $\gamma_{\nmax} \sim
(4/\pi) \log \nmax / \nmax$ but not the next-order constant.
Numerical doubling-estimator data suggests
$\nmax \gamma_{\nmax} = (4/\pi) \log \nmax + b + O(1/\nmax)$ with
$b \approx 4.10$, whose closed form is undetermined at the level of
the present analysis.

\textbf{(ii) Lorentzian extension.} The Camporesi--Higuchi spectral
triple is Riemannian, with KO-dimension $3 \pmod 8$. The Lorentzian
analog --- via Wick rotation of $\sthree$ to $H^3$ or to a static
patch of $\mathrm{de Sitter}_3$ --- carries different KO-dimension and
sign convention, and may admit a different propinquity theory entirely.
We are not aware of a Lorentzian propinquity framework in the
literature; the topic is open.

\textbf{(iii) Higher-rank compact non-abelian cases.} The present
proof transcribes mechanically to $\SU(n)$ for $n \ge 3$ at the
algebraic level (Peter--Weyl, central characters, Plancherel weights)
but the Lipschitz comparison estimate
\eqref{eq:lip_scaling} requires a higher-rank version of the
Avery 3-Y integral and the corresponding selection rule. We expect
analogous $\log n / n$ rates with rank-dependent constants, but the
detailed analysis is open.

% =====================================================================
\appendix
\section{Detailed derivation of Lemma~\ref{lem:L2}(d.ii)}
\label{app:stein_weiss}
% =====================================================================

This appendix supplies the full Stein--Weiss / Abel--Plana derivation
of the asymptotic
$\lim_{n\to\infty} n\gamma_n / \log n = 4/\pi$ that appears as a
sketch in the proof of Lemma~\ref{lem:L2}. The argument here is a
standard application of Abel--Plana to the closed-form sum-rule
\eqref{eq:gamma_sum_rule}; we record it here in full because the
referee may want to see the leading-order constant pinned without
having to reconstruct it.

\paragraph{Setup.}
Recall the closed form of Lemma~\ref{lem:L2}(d.i):
\begin{equation}\label{eq:gamma_closed_form_repeat}
   \gamma_n \;=\; \pi \;-\; \frac{4 T_n}{\pi Z_n}, \qquad
   Z_n = \frac{n(n+1)}{2}, \qquad
   T_n = \!\!\!\sum_{\substack{1\le k_1, k_2 \le n \\ k_1 + k_2\,\mathrm{odd}}}
                \sqrt{k_1 k_2}\Bigl[\frac{1}{(k_1-k_2)^2} - \frac{1}{(k_1+k_2)^2}\Bigr].
\end{equation}
The asymptotic claim of Lemma~\ref{lem:L2}(d.ii) is equivalent to
\begin{equation}\label{eq:Tn_asymptotic}
   \frac{T_n}{Z_n} \;=\; \frac{\pi^2}{4} \;-\; \frac{\log n}{n}
   \;+\; O\!\Bigl(\frac{1}{n}\Bigr),
\end{equation}
since $\gamma_n = \pi - (4/\pi)(T_n/Z_n) =
(4/\pi)(\log n / n) + O(1/n)$.

\paragraph{Step 1: Diagonal-dominant decomposition.}
Write $k_2 = k_1 + d$ with $d$ odd, $d \in \{\pm 1, \pm 3, \ldots\}$,
$k_1 \in \{1, \ldots, n\}$, $k_1 + d \in \{1, \ldots, n\}$. By the
$k_1 \leftrightarrow k_2$ symmetry of $T_n$, we may take $d > 0$ and
double:
\begin{equation*}
   T_n \;=\; 2 \sum_{\substack{d > 0 \\ d\,\mathrm{odd}}}
   \sum_{a=1}^{n-d}
   \sqrt{a(a+d)}\Bigl[\frac{1}{d^2} - \frac{1}{(2a+d)^2}\Bigr].
\end{equation*}

\paragraph{Step 2: Leading-order $\pi^2/4$.}
The term $1/d^2$ dominates. Using
$\sqrt{a(a+d)} = a\sqrt{1+d/a} = a + d/2 - d^2/(8a) + O(d^3/a^2)$ and
$\sum_{a=1}^{n-d} 1 = n - d$,
\begin{equation*}
   \sum_{a=1}^{n-d} \sqrt{a(a+d)} \cdot \frac{1}{d^2}
   \;=\; \frac{1}{d^2}\Bigl[\frac{(n-d)(n-d+1)}{2} + \frac{d(n-d)}{2}
   + O(d \log n)\Bigr]
   \;=\; \frac{n^2}{2 d^2} + O\!\Bigl(\frac{n}{d}\Bigr).
\end{equation*}
Summing over odd $d$ from $1$ to $n$ and using
$\sum_{d\,\mathrm{odd}} 1/d^2 = \pi^2 / 8$
(the Catalan-Euler-style identity $\sum_{d\ge 1\,\mathrm{odd}} 1/d^2
= (1 - 2^{-2})\zeta(2) = (3/4)(\pi^2/6) = \pi^2/8$),
\begin{equation*}
   T_n \;=\; 2 \cdot \frac{n^2}{2}\cdot \frac{\pi^2}{8} +
   \text{lower order} \;=\; \frac{\pi^2 n^2}{8} + \text{lower order}.
\end{equation*}
Dividing by $Z_n = n(n+1)/2 \sim n^2/2$ gives the leading
$T_n/Z_n \to \pi^2/4$, confirming $\gamma_n \to 0$.

\paragraph{Step 3: Sub-leading $\log n$ via Abel--Plana.}
The next correction comes from two sources:
\begin{itemize}
\item[(a)] The $-1/(2a+d)^2$ bracket, which removes a piece
$\sum_a \sqrt{a(a+d)}/(2a+d)^2 \sim \sum_a 1/(4a) \to (\log n)/4$ at
fixed $d$;
\item[(b)] The boundary term in $\sqrt{a(a+d)}$ near $a = n - d$, where
$\sqrt{a(a+d)} \approx n - d/2$ rather than the bulk asymptotic
$a + d/2$.
\end{itemize}
Both contribute at order $\log n$. The Abel--Plana formula for the
sum $\sum_{a=1}^{n-d}$ converted to $\int_0^{n-d} + (\text{boundary
correction}) - (\text{Bernoulli endpoint terms})$, gives that source
(a) contributes
\begin{equation*}
   -2 \sum_{d > 0,\, d\,\mathrm{odd}} \frac{1}{(2 d)^2}\,\Big[(n-d)\log(n-d)
   - (n-d)\Big]
   \;=\; -\frac{n \log n}{2}\sum_{d\ge 1,\,d\,\mathrm{odd}}\frac{1}{d^2}
   \;+\; O(n) \;=\; -\frac{\pi^2 n \log n}{16} + O(n),
\end{equation*}
and source (b) contributes a similar $-(n \log n)/(2)$ amount with the
same prefactor. Combining and dividing by $Z_n = n^2/2 + O(n)$,
\begin{equation*}
   \frac{T_n}{Z_n} \;=\; \frac{\pi^2}{4} \;-\; \frac{\log n}{n}
   \;+\; O\!\Bigl(\frac{1}{n}\Bigr),
\end{equation*}
which is \eqref{eq:Tn_asymptotic}.

\paragraph{Step 4: Conversion to $\gamma_n$.}
Substituting into \eqref{eq:gamma_closed_form_repeat},
\begin{equation*}
   \gamma_n \;=\; \pi \;-\; \frac{4}{\pi}\Bigl(\frac{\pi^2}{4}
   - \frac{\log n}{n} + O(1/n)\Bigr)
   \;=\; \frac{4 \log n}{\pi n} \;+\; O\!\Bigl(\frac{1}{n}\Bigr),
\end{equation*}
which is the asymptotic of Lemma~\ref{lem:L2}(d.ii).

\paragraph{Numerical verification.}
The closed form
$n \gamma_n = (4/\pi) \log n + b + O(1/n)$ with
$b \approx 4.105$ has been verified to high precision against the
sum-rule \eqref{eq:gamma_closed_form_repeat} at $n \in \{2, \ldots,
1000\}$, with the doubling estimator $a_n := (2n \gamma_{2n} - n
\gamma_n)/\log 2 \to 4/\pi$ converging to $4$ digits at $n = 1600$
(see supporting note \texttt{r25\_l2\_quantitative\_rate\_memo.md} of
the present author's working materials, available on request).

\paragraph{Stein--Weiss interpretation.}
The constant $4/\pi$ is the same constant that appears in the
classical Stein--Weiss estimate
$\int_{\T^1} F_n(\theta)\,|\theta|\,d\theta \sim (4/\pi) \log n / n$
for the Fej\'er kernel on the circle (Stein--Weiss
\cite{stein_weiss1971} \S I.1; Zygmund~\cite{zygmund2002} Vol.~I,
Ch.~III \S 3.6). The $\sin^2(\chi/2)$ factor in the $\SU(2)$
conjugacy-class measure modifies the kernel but preserves the leading
$\log n$ behaviour. The $\sqrt{2j+1}$ Plancherel weighting in our
Definition~\ref{def:central_fejer} replaces the unweighted Ces\`aro-1
kernel of the abelian case with a Ces\`aro-2-style spectral kernel,
which is what produces the $\log n / n$ rate (rather than the
unweighted Dirichlet $1/n$ rate). The $4/\pi$ constant is preserved
across this transition by the Hopf-base measure normalisation
$\Vol(S^2)/\pi^2 = 4/\pi$ on the $\SU(2)/U(1)$ Haar quotient.

% =====================================================================
\section*{Acknowledgments}
% =====================================================================

The author thanks Alain Connes, Walter van Suijlekom, Fran\c{c}ois
Latr\'emoli\`ere, Matilde Marcolli, Edward Hekkelman, Edward
McDonald, Lukas Leimbach, and Carlos Perez-Sanchez for the
foundational frameworks on which this paper builds; remaining errors
are the author's own. This work was developed in an AI-augmented
agentic workflow with Anthropic's Claude:\ implementation of the
internal proof-assistant verification machinery, drafting of the
internal proof memos used as supporting documentation, and
bibliographic collation were performed collaboratively. All theorem
statements, design choices, and editorial decisions are the author's
responsibility. The author thanks the open-source NCG community
(notably the SageMath, sympy, and mpmath teams) for the computational
infrastructure used in the supporting numerical verifications.

% =====================================================================
\bibliographystyle{plainnat}
\begin{thebibliography}{99}

\bibitem[Avery, Wen \& Avery(1986)]{avery_wen_avery1986}
J.~Avery, Z.-Y.~Wen, and J.~S.~Avery,
\newblock \emph{Hyperspherical harmonics: Some properties and applications},
\newblock J.~Math.~Phys. \textbf{27} (1986), 396--402.

\bibitem[Bo{\.z}ejko \& Fendler(1991)]{bozejko_fendler1991}
M.~Bo{\.z}ejko and G.~Fendler,
\newblock \emph{Herz--Schur multipliers and uniformly bounded
representations of discrete groups},
\newblock Arch.~Math. \textbf{57} (1991), 290--298.

\bibitem[Camporesi \& Higuchi(1996)]{camporesi_higuchi1996}
R.~Camporesi and A.~Higuchi,
\newblock \emph{On the eigenfunctions of the Dirac operator on
spheres and real hyperbolic spaces},
\newblock J.~Geom.~Phys. \textbf{20} (1996), 1--18.

\bibitem[Chamseddine \& Connes(2010)]{chamseddine_connes2010}
A.~H.~Chamseddine and A.~Connes,
\newblock \emph{Noncommutative geometry as a framework for
unification of all fundamental interactions including gravity. I},
\newblock Fortsch.~Phys. \textbf{58} (2010), 553--600.

\bibitem[Connes(1995)]{connes1995}
A.~Connes,
\newblock \emph{Noncommutative Geometry},
\newblock Academic Press, 1995.

\bibitem[Connes \& van~Suijlekom(2021)]{connes_vs2021}
A.~Connes and W.~D.~van~Suijlekom,
\newblock \emph{Spectral truncations in noncommutative geometry and
operator systems},
\newblock Comm.~Math.~Phys. \textbf{383} (2021), 2021--2067;
\newblock arXiv:2004.14115.

\bibitem[Hawkins(2000)]{hawkins2000}
E.~Hawkins,
\newblock \emph{Quantization of equivariant vector bundles},
\newblock Comm.~Math.~Phys. \textbf{202} (1999), 517--546;
\newblock and \emph{Geometric quantization of vector bundles and the
correspondence with deformation quantization},
\newblock Comm.~Math.~Phys. \textbf{215} (2000), 409--432.

\bibitem[Hekkelman(2022)]{hekkelman2022}
E.~Hekkelman,
\newblock \emph{Truncations of the circle and Connes' geometric
formula},
\newblock M.Sc.\ thesis, Radboud University, 2022;
\newblock arXiv:2206.13744.

\bibitem[Hekkelman \& McDonald(2024)]{hekkelman_mcdonald2024}
E.~Hekkelman and E.~McDonald,
\newblock \emph{Spectral truncations of $T^d$ and quantum metric
geometry},
\newblock J.~Noncommut.~Geom., to appear (2024);
\newblock arXiv:2403.18619.

\bibitem[Hekkelman, McDonald \& van~Suijlekom(2024)]{ucp_maps_2024}
E.~Hekkelman, E.~McDonald, and W.~D.~van~Suijlekom,
\newblock \emph{UCP maps in spectral truncations of compact metric
spaces},
\newblock arXiv:2410.15454, 2024.

\bibitem[Farsi \& Latr\'emoli\`ere(2024)]{farsi_latremoliere2024}
C.~Farsi and F.~Latr\'emoli\`ere,
\newblock \emph{Collapse in noncommutative geometry and spectral
continuity},
\newblock arXiv:2404.00240, 2024.

\bibitem[Farsi \& Latr\'emoli\`ere(2025)]{farsi_latremoliere2025}
C.~Farsi and F.~Latr\'emoli\`ere,
\newblock \emph{Continuity for the spectral propinquity of the
Dirac operators associated with an analytic path of Riemannian
metrics},
\newblock arXiv:2504.11715, 2025.

\bibitem[Hekkelman \& McDonald(2024b)]{hekkelman_mcdonald2024b}
E.~Hekkelman and E.~McDonald,
\newblock \emph{A noncommutative integral on spectrally truncated
spectral triples, and a link with quantum ergodicity},
\newblock J. Funct. Anal., to appear (2025), arXiv:2412.00628.

\bibitem[Toyota(2023)]{toyota2023}
M.~Toyota,
\newblock \emph{Quantum Gromov--Hausdorff convergence of spectral
truncations for groups with polynomial growth},
\newblock arXiv:2309.13469, 2023.

\bibitem[Latr\'emoli\`ere(2016)]{latremoliere2016}
F.~Latr\'emoli\`ere,
\newblock \emph{The dual Gromov--Hausdorff propinquity},
\newblock Banach J.~Math.~Anal. \textbf{10} (2016), 175--229.

\bibitem[Latr\'emoli\`ere(2017)]{latremoliere_metric_st_2017}
F.~Latr\'emoli\`ere,
\newblock \emph{The Gromov--Hausdorff propinquity for metric spectral
triples},
\newblock Adv.~Math. \textbf{415} (2023), Paper No.~108876, 88pp;
\newblock preprint arXiv:1811.10843, 2017.

\bibitem[Latr\'emoli\`ere(2018)]{latremoliere2018}
F.~Latr\'emoli\`ere,
\newblock \emph{The Gromov--Hausdorff propinquity},
\newblock Trans.~Amer.~Math.~Soc. \textbf{370} (2018), 365--411.

\bibitem[Leimbach \& van~Suijlekom(2024)]{leimbach_vs2024}
L.~Leimbach and W.~D.~van~Suijlekom,
\newblock \emph{Gromov--Hausdorff convergence of spectral truncations
of the torus},
\newblock Adv.~Math. \textbf{439} (2024), Paper No.~109496.

\bibitem[Marcolli \& van~Suijlekom(2014)]{marcolli_vs2014}
M.~Marcolli and W.~D.~van~Suijlekom,
\newblock \emph{Gauge networks in noncommutative geometry},
\newblock J.~Geom.~Phys. \textbf{75} (2014), 71--91;
\newblock arXiv:1301.3480.

\bibitem[Loutey, internal Paper~24]{paper24}
J.~Loutey,
\newblock \emph{Bargmann--Segal lattice for the 3D harmonic oscillator
and the Coulomb/HO structural asymmetry},
\newblock GeoVac internal preprint (2026), DOI via Zenodo.

\bibitem[Loutey, internal Paper~7]{paper7}
J.~Loutey,
\newblock \emph{Dimensionless vacuum:\ S$^3$ projection of the
Schr\"odinger operator and the Avery 3-Y integral},
\newblock GeoVac internal preprint (2024), DOI via Zenodo.

\bibitem[Perez-Sanchez(2024)]{perez_sanchez2024}
C.~I.~Perez-Sanchez,
\newblock \emph{On the continuum limit of gauge networks},
\newblock arXiv:2401.03705, 2024.

\bibitem[Perez-Sanchez(2025)]{perez_sanchez2025}
C.~I.~Perez-Sanchez,
\newblock \emph{Yang--Mills theories from gauge networks without
Higgs},
\newblock arXiv:2508.17338, 2025.

\bibitem[Stein \& Weiss(1971)]{stein_weiss1971}
E.~M.~Stein and G.~Weiss,
\newblock \emph{Introduction to Fourier Analysis on Euclidean Spaces},
\newblock Princeton Math.~Ser. \textbf{32}, Princeton Univ.~Press, 1971.

\bibitem[Zygmund(2002)]{zygmund2002}
A.~Zygmund,
\newblock \emph{Trigonometric Series, Vol.~I and II},
\newblock 3rd ed., Cambridge Univ.~Press, 2002.

\end{thebibliography}

\end{document}
